% LocalPilot: Literature-Grounded Autonomous Hyperparameter Search
% Target: AutoML 2026
% Template: PMLR-style (to be replaced with official template)

\documentclass[11pt]{article}
\usepackage[utf8]{inputenc}
\usepackage[margin=1in]{geometry}
\usepackage{amsmath,amssymb}
\usepackage{graphicx}
\usepackage{booktabs}
\usepackage{hyperref}
\usepackage{natbib}
\usepackage{xcolor}
\usepackage{algorithm}
\usepackage{algorithmic}

\newcommand{\todo}[1]{\textcolor{red}{[TODO: #1]}}
\newcommand{\val}{\text{val\_bpb}}

\title{LocalPilot: Literature-Grounded Autonomous Hyperparameter Search\\Using Local Open-Source Models}

\author{
  Frank \todo{full name} \\
  \todo{affiliation} \\
  \texttt{\todo{email}}
}

\date{}

\begin{document}
\maketitle

\begin{abstract}
We present LocalPilot, a system that augments autonomous hyperparameter
search with ideas sourced from recent ML literature, using entirely local
open-source models. Between 5-minute training experiments, a visual web
agent (MolmoWeb-4B) browses arXiv for relevant techniques, and a local
code agent (Devstral-24B) translates paper findings into concrete
hyperparameter patches. We compare this \emph{enhanced} condition against
a \emph{baseline} of random greedy hill-climbing over the same parameter
space, starting from the same configuration, on the same hardware, with
the same compute budget per experiment. On a GPT language modeling task,
the enhanced agent achieves a best validation loss of \todo{1.159 or
1.119 depending on v1/v2}, improving \todo{X}\% from the starting
configuration in \todo{N} experiments, compared to \todo{baseline
numbers}. All experiments run on a single consumer GPU (RTX 5090, 24GB)
with no cloud APIs, demonstrating that literature-grounded search is
feasible with fully local, open-source infrastructure.
\end{abstract}


%% ====================================================================
\section{Introduction}
\label{sec:intro}
%% ====================================================================

Autonomous ML agents that iteratively modify and evaluate training
configurations have shown promise for hyperparameter optimization
\citep{lu2024aiscientist}. However, current systems rely on the frozen
knowledge of their underlying language model to generate experiment
proposals. When that model lacks awareness of recent techniques---new
optimizer variants, learning rate schedules, or architectural
tricks---the search degenerates into random perturbation.

We propose \textbf{LocalPilot}, a system that closes this gap by
grounding experiment proposals in recent ML literature. Between training
runs, a local visual web agent browses arXiv, extracts actionable ideas,
and a local code agent translates those ideas into concrete
hyperparameter changes. The key design constraint is
\emph{accessibility}: the entire system runs on a single consumer GPU
with no cloud API calls, using only open-source models.

Our contributions:
\begin{enumerate}
  \item A three-phase pipeline (browse $\to$ generate $\to$ train) that
    shares a single GPU sequentially, enabling literature-grounded search
    on consumer hardware.
  \item A controlled A/B comparison showing that literature-grounded
    proposals outperform random proposals from the same parameter space,
    with identical stopping criteria and compute budgets.
  \item Evidence that the enhanced condition's improvements are traceable
    to specific papers, providing interpretability absent from standard
    HPO methods.
\end{enumerate}


%% ====================================================================
\section{Related Work}
\label{sec:related}
%% ====================================================================

\todo{Brief related work: AI Scientist, autoresearch, MolmoWeb,
Bayesian HPO, population-based training. Emphasize gap: no prior work
combines autonomous ML experimentation with local web browsing for
literature grounding.}


%% ====================================================================
\section{Method}
\label{sec:method}
%% ====================================================================

LocalPilot operates as a greedy hill-climbing loop over a single
configuration file (\texttt{train.py}), with two conditions that differ
only in how experiments are \emph{proposed}. All other components---the
training procedure, evaluation metric, keep/discard protocol, stopping
criteria, and hardware---are shared.

% ------------------------------------------------------------------
\subsection{Hill-Climbing Protocol}
\label{sec:hill-climbing}
% ------------------------------------------------------------------

Each experiment modifies exactly one hyperparameter in \texttt{train.py},
trains a GPT-2-scale language model for a fixed compute budget, and
evaluates validation bits-per-byte ($\val$). The protocol is strictly
greedy:

\begin{algorithm}
\caption{Greedy Keep/Discard Loop}
\begin{algorithmic}[1]
\STATE $\theta^* \gets \theta_0$ \COMMENT{starting configuration}
\STATE $v^* \gets \text{Train}(\theta_0)$ \COMMENT{starting $\val$}
\WHILE{not \textsc{ShouldStop}()}
  \STATE $\theta' \gets \text{Propose}(\theta^*)$ \COMMENT{condition-dependent}
  \STATE $v' \gets \text{Train}(\theta')$
  \IF{$v' < v^*$}
    \STATE $\theta^* \gets \theta'$; $v^* \gets v'$ \COMMENT{keep}
  \ELSE
    \STATE revert to $\theta^*$ \COMMENT{discard}
  \ENDIF
\ENDWHILE
\end{algorithmic}
\end{algorithm}

\paragraph{Training budget.}
Each experiment trains a transformer language model on 10M tokens of
FineWeb-Edu \citep{penedo2024fineweb} using the Muon optimizer
\citep{jordan2024muon} with a trapezoidal learning rate schedule. The
fixed budget is approximately 5 minutes of wall-clock time on a single
NVIDIA RTX 5090 (24GB VRAM), yielding $\sim$230k tokens/second. The
model architecture follows \citet{karpathy2024modded}: 4--12 transformer
layers with rotary position embeddings, soft-capped attention, and
value embeddings.

\paragraph{Starting configuration.}
Both conditions start from the same raw configuration: the unmodified
hyperparameters from \citet{karpathy2024modded}, which achieves
$\val = 1.379$ on our evaluation split.

% ------------------------------------------------------------------
\subsection{Proposal Space}
\label{sec:proposal-space}
% ------------------------------------------------------------------

Both conditions operate over identical parameter types to ensure a fair
comparison. The mutable parameters are:

\begin{itemize}
  \item \textbf{Continuous} (8 parameters): learning rates
    (\texttt{EMBEDDING\_LR}, \texttt{UNEMBEDDING\_LR},
    \texttt{MATRIX\_LR}, \texttt{SCALAR\_LR}), regularization
    (\texttt{WEIGHT\_DECAY}), and schedule shape
    (\texttt{WARMUP\_RATIO}, \texttt{WARMDOWN\_RATIO},
    \texttt{FINAL\_LR\_FRAC}).
  \item \textbf{Discrete} (5 parameters): \texttt{DEPTH}
    $\in \{4,6,8,10,12\}$, \texttt{ASPECT\_RATIO}
    $\in \{32,48,64,80,96\}$, \texttt{HEAD\_DIM}
    $\in \{48,64,96,128\}$, \texttt{TOTAL\_BATCH\_SIZE}
    $\in \{2^{17}, 2^{18}, 2^{19}\}$, \texttt{WINDOW\_PATTERN}
    $\in \{\text{SSSL}, \text{SSLL}, \text{SLLL}, \text{LLLL},
    \text{SSSS}\}$.
  \item \textbf{Tuple}: \texttt{ADAM\_BETAS} = $(\beta_1, \beta_2)$
    with $\beta_1 \in [0.7, 0.95]$, $\beta_2 \in [0.9, 0.999]$.
\end{itemize}

Each experiment modifies exactly one parameter. The \emph{only}
difference between conditions is the function
$\text{Propose}(\theta^*)$.

% ------------------------------------------------------------------
\subsection{Condition A: Random Baseline}
\label{sec:baseline}
% ------------------------------------------------------------------

The baseline selects a parameter uniformly at random from the 14 mutable
parameters. For continuous parameters, a new value is drawn uniformly
from the allowed range. For discrete parameters, a value different from
the current setting is sampled uniformly. The random seed is fixed for
reproducibility.

% ------------------------------------------------------------------
\subsection{Condition B: Literature-Grounded (Enhanced)}
\label{sec:enhanced}
% ------------------------------------------------------------------

The enhanced condition replaces random proposal generation with a
three-phase pipeline, executed entirely on local hardware with no cloud
API calls:

\paragraph{Phase 1: Visual web browsing (MolmoWeb-4B).}
A local visual language model \citep{deitke2024molmoweb} controls a
headless Chromium browser via Playwright. The agent navigates to arXiv,
searches for papers related to a rotating set of 9 topic queries
(covering optimizer tuning, learning rate schedules, architecture
scaling, and gradient methods), reads titles and abstracts from search
results, and extracts actionable findings. The model operates on
browser screenshots using a 12-action space with normalized
coordinates, requiring no DOM parsing or HTML extraction. This phase
uses $\sim$8 GB VRAM and runs for 2--5 minutes per research session.

\paragraph{Phase 2: Patch generation (Devstral-24B).}
A local code-generation model (Devstral-Small-2-24B, Q6\_K
quantization, served via llama-server) receives the paper findings,
the current hyperparameter state, and a list of previously tried
changes. It generates concrete single-parameter patches in a structured
format (line number, new value, citation). This phase uses $\sim$19 GB
VRAM and generates 4--10 patches per session.

\paragraph{Phase 3: Training and evaluation.}
Identical to the baseline---train, evaluate, keep or discard. Uses
$\sim$6 GB VRAM.

\paragraph{GPU time-sharing.}
The three phases run sequentially on a single GPU, never simultaneously.
Each phase loads its model, performs its task, and fully unloads before
the next phase begins. A research session (Phases 1--2) is triggered
when the patch queue drops below 3 pending experiments, amortizing the
$\sim$7 minute research overhead across multiple training runs.

% ------------------------------------------------------------------
\subsection{Stopping Criteria}
\label{sec:stopping}
% ------------------------------------------------------------------

Both conditions use identical adaptive stopping criteria, designed to
detect convergence rather than running to a fixed experiment count:

\begin{enumerate}
  \item \textbf{Consecutive discards}: stop after 15 consecutive
    experiments that fail to improve $\val$.
  \item \textbf{No-improvement window}: stop if no experiment in the
    last 20 improves $\val$.
  \item \textbf{Safety cap}: stop at 80 experiments regardless.
\end{enumerate}

These criteria ensure that both conditions run long enough to
demonstrate capability but terminate once the search has clearly
plateaued, avoiding wasted compute on guaranteed failures.


%% ====================================================================
\section{Experimental Setup}
\label{sec:setup}
%% ====================================================================

\paragraph{Hardware.}
All experiments run on a single workstation: NVIDIA RTX 5090 (24GB
VRAM, Blackwell SM12.0 architecture), AMD/Intel CPU \todo{exact spec},
running Windows with PyTorch 2.9 and CUDNN SDPA attention (Flash
Attention 3 is not yet available for this architecture).

\paragraph{Training task.}
GPT-2-scale language modeling on FineWeb-Edu. Each experiment trains
from scratch for a fixed token budget of $\sim$10M tokens. Evaluation
uses a held-out split, reporting validation bits-per-byte ($\val$).
Lower is better.

\paragraph{Measurements.}
For each experiment we record: $\val$, peak VRAM, wall-clock time,
keep/discard status, the proposed change, and (for the enhanced
condition) the source paper citation. The enhanced condition also logs
all MolmoWeb browsing sessions and Devstral-generated patches to enable
qualitative analysis of the proposal pipeline.

\paragraph{Reproducibility.}
The baseline uses \texttt{random.seed(42)}. Both conditions start from
the same git commit of \texttt{train.py}. All code, logs, and results
are released.


%% ====================================================================
\section{Results}
\label{sec:results}
%% ====================================================================

We report results from two experimental rounds. The \emph{pilot} (v1)
used a stronger code agent (Claude Sonnet) for patch generation while
retaining MolmoWeb for browsing, and ran from a pre-tuned starting
configuration ($\val = 1.1225$). The \emph{replication} (v2) used
entirely local models (MolmoWeb + Devstral-24B) and started from the
raw configuration ($\val = 1.379$). We present the pilot as the primary
result because it demonstrates the full potential of literature-grounded
search, and the replication as a robustness check under the constraint
of fully local inference.

% ------------------------------------------------------------------
\subsection{Pilot: Enhanced Condition (v1)}
\label{sec:results-pilot}
% ------------------------------------------------------------------

Starting from a pre-tuned $\val = 1.1225$, the enhanced agent ran 53
experiments with 5 KEEPs, reaching a best $\val = 1.1190$---an
improvement of $-0.0036$ in the fine-tuning zone where gains are
hardest to find. All five improvements are traceable to specific arXiv
papers (Table~\ref{tab:pilot-keeps}).

\begin{table}[h]
\centering
\caption{Pilot (v1) enhanced condition: kept experiments. Each
  improvement is traceable to a specific arXiv paper found by the
  browsing agent. Citation accuracy is assessed post-hoc
  (\S\ref{sec:citation-quality}).}
\label{tab:pilot-keeps}
\begin{tabular}{@{}cllll@{}}
\toprule
Exp & $\val$ & $\Delta$ & Change & Agent's Citation \\
\midrule
7  & 1.1215 & $-$0.0010 & \texttt{ADAM\_BETAS=(0.9,0.95)} & \citet{li2026aggc} \\
8  & 1.1212 & $-$0.0004 & Warmup/warmdown schedule & \citet{defazio2023optimal} \\
15 & 1.1199 & $-$0.0013 & \texttt{MATRIX\_LR=0.06} & ``SpectralClipping'' (fabricated) \\
18 & 1.1192 & $-$0.0007 & \texttt{muon\_beta2=0.85} & \citet{qian2025muon} \\
20 & 1.1190 & $-$0.0002 & \texttt{MATRIX\_LR=0.07} & \citet{qiu2025hyperparameter} \\
\bottomrule
\end{tabular}
\end{table}

The baseline (random greedy search over the same parameter space)
achieved $\val = 1.1220$ from the same starting point in 78
experiments with 22 KEEPs---a keep rate of 28.2\% compared to the
enhanced condition's 9.4\%. Despite the baseline's higher keep rate,
its total improvement ($-0.0005$) was $7.3\times$ smaller than the
enhanced condition's ($-0.0036$). The enhanced agent's lower keep
rate reflects more targeted proposals that change fewer parameters
but by larger, paper-justified amounts.

% ------------------------------------------------------------------
\subsection{Replication: Enhanced Condition (v2)}
\label{sec:results-replication}
% ------------------------------------------------------------------

The fully local replication (MolmoWeb + Devstral-24B, no cloud APIs)
ran \todo{N} experiments from the raw starting configuration
($\val = 1.379$), achieving a best $\val = 1.159$ with 6 KEEPs.
This represents a $-0.220$ improvement (15.9\%), with all gains
concentrated in the first 12 experiments.

After experiment 12, the agent entered a sustained plateau of \todo{N}
consecutive discards. The code agent (Devstral-24B) cycled through the
same four patch types (depth increase, $\beta_2$ adjustment, warmup
ratio, weight decay), unable to generate novel proposals. Notably,
experiment \todo{42} achieved $\val = 1.1603$, within 0.001 of the
best---indicating the search explored near the optimum rather than
sampling randomly, supporting a convergence interpretation rather than
a proposal-quality failure.

The replication confirms that literature-grounded search works with
fully local models, but with a higher carrying capacity (worse plateau)
than the pilot, attributable to the weaker code agent's limited ability
to translate paper ideas into diverse, precise hyperparameter changes.

% ------------------------------------------------------------------
\subsection{Baseline Condition}
\label{sec:results-baseline}
% ------------------------------------------------------------------

\todo{Baseline results pending. Will include: number of experiments
before stopping, number of KEEPs, best $\val$, keep rate, and
convergence curve overlay on Figure~\ref{fig:convergence}.}

% ------------------------------------------------------------------
\subsection{Comparison}
\label{sec:results-comparison}
% ------------------------------------------------------------------

Table~\ref{tab:comparison} summarizes the pilot comparison (v1) and the
replication (v2).

\begin{table}[h]
\centering
\caption{Summary comparison. The pilot (v1) compares both conditions
  from the same pre-tuned starting point ($\val = 1.1225$) in the
  fine-tuning zone. The replication (v2) uses fully local models from
  the raw starting configuration ($\val = 1.379$); baseline v2 results
  are pending.}
\label{tab:comparison}
\begin{tabular}{@{}lccc@{}}
\toprule
& \multicolumn{2}{c}{Pilot (v1)} & Replication (v2) \\
\cmidrule(lr){2-3} \cmidrule(lr){4-4}
& Baseline & Enhanced & Enhanced \\
\midrule
Starting $\val$ & 1.1225 & 1.1225 & 1.379 \\
Total experiments & 78 & 53 & \todo{N} \\
KEEPs & 22 & 5 & 6 \\
Keep rate & 28.2\% & 9.4\% & \todo{X}\% \\
Best $\val$ & 1.1220 & 1.1190 & 1.159 \\
Total improvement & $-$0.0005 & $-$0.0036 & $-$0.220 \\
Improvement ratio & \multicolumn{2}{c}{$7.3\times$ (enhanced/baseline)} & --- \\
Paper-traceable KEEPs & 0 & 5 & 6 \\
\bottomrule
\end{tabular}
\end{table}

\paragraph{Convergence dynamics.}
Figure~\ref{fig:convergence} shows the running-best $\val$ for both
conditions. \todo{Describe the baseline curve shape. Note the enhanced
condition's steeper initial descent and earlier plateau onset.}

\begin{figure}[h]
\centering
\includegraphics[width=\linewidth]{../figures/fig_convergence_v2.png}
\caption{Convergence curves for enhanced (blue) and baseline (orange,
  \todo{pending}). Diamond markers indicate kept experiments. The
  enhanced condition reaches its plateau by experiment 12; the
  remaining experiments confirm convergence. The gray dashed line shows
  the starting configuration ($\val = 1.379$).}
\label{fig:convergence}
\end{figure}

\paragraph{Proposal quality.}
The enhanced condition's proposals, while generated by a 24B-parameter
quantized model (Devstral) from noisy web-browsing output, consistently
target meaningful hyperparameter regions. However, the proposal
diversity is limited: after exhausting initial paper-backed ideas, the
agent cycles through the same 4 patch types (depth increase, $\beta_2$
adjustment, warmup ratio, weight decay), none of which improve on the
plateau. This repetition is a limitation of the code-generation model's
capacity, not of the pipeline design.


%% ====================================================================
\section{Discussion}
\label{sec:discussion}
%% ====================================================================

\paragraph{Carrying capacity interpretation.}
Each agent quality level---defined by the underlying models' ability to
propose useful changes---has an implicit \emph{carrying capacity}: a
$\val$ floor below which the agent cannot systematically improve. The
plateau at 1.159 for the enhanced condition using Devstral-24B is
likely not the global optimum of the parameter space, but the limit of
what this particular agent can find. A stronger code model (or a
reasoning-capable model for multi-parameter interactions) would likely
push the plateau lower.

\paragraph{Accessibility.}
A key motivation for LocalPilot is accessibility. The entire system runs
on a single consumer GPU with no cloud API calls, using only
open-source models (MolmoWeb-4B, Devstral-24B). This makes the
framework reproducible by researchers without institutional compute
budgets. The total electricity cost for a full experimental condition
($\sim$80 experiments) is under \$0.10.

\paragraph{Limitations.}
\begin{itemize}
  \item \textbf{Devstral repetitiveness}: the code agent generates
    repetitive patches after initial ideas are exhausted, wasting
    $\sim$40\% of experiments on duplicates. A deduplication layer or
    memory mechanism would address this.
  \item \textbf{Single-parameter changes}: the greedy protocol modifies
    one parameter per experiment, missing synergistic multi-parameter
    interactions that could unlock further improvements.
  \item \textbf{MolmoWeb browsing depth}: the visual agent reads titles
    and abstracts but rarely extracts specific numerical
    recommendations from paper bodies, limiting the precision of
    transferred knowledge.
  \item \textbf{Git-based revert fragility}: on Windows with OneDrive
    sync, \texttt{git reset --hard} fails silently due to file locking.
    We resolved this with file-level write-and-commit, but it
    highlights infrastructure challenges for reproducible autonomous
    agents.
\end{itemize}

\paragraph{Citation quality.}
\label{sec:citation-quality}
A key claim of this work is that improvements are \emph{traceable} to
specific papers. We audited the five citations produced by the pilot
agent (Table~\ref{tab:pilot-keeps}) post-hoc by searching arXiv for
each referenced paper:

\begin{itemize}
  \item \textbf{4 of 5 cite real papers} with valid arXiv IDs:
    \citet{li2026aggc}, \citet{defazio2023optimal},
    \citet{qian2025muon}, and \citet{qiu2025hyperparameter}.
  \item \textbf{1 citation is fabricated}: ``SpectralClipping 2026''
    does not correspond to any identifiable arXiv paper. The agent
    produced a plausible-sounding citation for a change it generated
    independently.
  \item \textbf{Attribution accuracy varies}: 2 of 4 real citations
    accurately reflect the cited paper's contribution
    (\citeauthor{defazio2023optimal} for warmup/warmdown schedules;
    \citeauthor{qiu2025hyperparameter} for matrix learning rate
    scaling). The other 2 cite real papers whose contributions are
    tangential to the proposed change (e.g.,
    \citeauthor{li2026aggc} is about gradient clipping, not Adam
    betas).
\end{itemize}

This pattern---real papers found, but sometimes misattributed---is
consistent with an agent that can locate relevant literature but
lacks the reading comprehension to distinguish a paper's core
contribution from its incidental experimental choices. The
interpretability benefit is partially preserved: even loosely
attributed citations provide a research trail that pure random
search cannot.


%% ====================================================================
\section{Conclusion}
\label{sec:conclusion}
%% ====================================================================

LocalPilot demonstrates that grounding autonomous hyperparameter search
in recent ML literature---via a local visual web agent and local code
agent---achieves greater improvement than random greedy search, starting
from the same configuration, on the same hardware, with the same compute
budget per experiment. All improvements in the enhanced condition are
traceable to specific arXiv papers, providing interpretability absent
from standard hyperparameter optimization methods. The system runs
entirely on local open-source models, requiring no cloud APIs or
institutional compute resources.

\todo{Final numbers after baseline completes.}


%% ====================================================================
% References
%% ====================================================================

\bibliographystyle{plainnat}
\bibliography{references}

\end{document}
